\documentclass{article}

\usepackage[english]{babel}
\usepackage[letterpaper,top=2cm,bottom=2cm,left=3cm,right=3cm,marginparwidth=1.75cm]{geometry}
\usepackage{amsmath}
\usepackage{graphicx}
\usepackage{amssymb}
\usepackage{float}
\usepackage{yhmath}
\usepackage[colorlinks=true, allcolors=blue]{hyperref}

\DeclareMathOperator{\sgn}{sgn}
\newcommand*\conj[1]{\bar{#1}}
\newcommand*\mean[1]{\bar{#1}}
\newcommand{\RE}{\mathrm{Re}}
\newcommand{\IM}{\mathrm{Im}}
\counterwithin*{equation}{section}
\renewcommand\theequation{\thesection.\arabic{equation}}
\newcommand{\infint}{\displaystyle\int_{-\infty}^{+\infty}}

\title{Controls on the distribution of extensional strain between rifting events\\
  \large Supplementary Derivations}
\author{Navarrete, Iv\'an \and Olive, Jean-Arthur}

% ====================================================================
\begin{document}
\maketitle

% ====================================================================
\section{Introduction}

Continental rifting is a fundamental tectonic process that facilitates the
transition from stable cratonic lithosphere to the formation of new basins
(Corti, 2009; Ebinger et al., 2013). This evolution is characterized by
rifting cycles that operate across a wide range of spatio-temporal scales,
from rift initiation to incipient break-up (Ebinger, 1989; Le Gall et al.,
2008). These cycles are typically described in three distinct phases:
(i) the near-instantaneous, elastic strain release through dike emplacement
and normal faulting (Grandin et al., 2011; Maccaferri et al., 2016);
(ii) a period of post-diking visco-elastic relaxation (Foulger et al., 1992;
Nooner et al., 2009); and (iii) the inter-rifting phase, which reflects the
steady, secular divergence of the plates (Saria et al., 2014; Stamps et al.,
2021). While the first two phases are well-documented by geodetic observations
of modern rifting episodes such as those in Iceland (Heki et al., 1993) and
the Afar (Nooner et al., 2009), the mechanisms governing the build-up of
extensional strain \emph{between} events remain the least well-constrained
aspect of the rifting cycle.

\subsection{Limitations of Existing Kinematic Models}

The traditional framework for modeling inter-rifting strain relies on the
``backslip'' or ``locking'' models originally developed for the earthquake
cycle (Savage and Prescott, 1978). In this paradigm the rift is treated as a
dislocation within an elastic half-space that opens continuously at a
depth-dependent rate (LaFemina et al., 2005; Konca et al., 2018). These models
face significant physical challenges. The elastic half-space assumption ignores
the well-documented rheological transition at rifts, where high heat flow and
the presence of melt weaken the lower crust and upper mantle (Prodehl et al.,
1997; Plasman et al., 2017). An infinitely deep, continuously opening dike is
moreover mechanically untenable and inconsistent with the short visco-elastic
relaxation times expected in magmatically active regions (Ebinger et al.,
2013; Weinstein et al., 2017). At the numerical level, finite-domain
simulations must impose vanishing velocities at the boundaries, a condition
that cannot be reconciled with the constant far-field divergence observed in
global reference frames (Fernandes et al., 2004; Altamimi et al., 2016).
The role of magma further complicates the inter-rifting signal: volatile
release and the presence of melt fractions in the asthenosphere (Lee et al.,
2016; Muirhead et al., 2016; Holtzman, 2016) evolve the mechanical strength of
the lithosphere on timescales that overlap with the rifting cycle itself.
Consequently, a ``kinematic gap'' exists between the transient decay of a
single rifting event and the steady-state strain rates predicted by global
plate motion models (Kreemer et al., 2014; Stamps et al., 2018).

\subsection{A Unified Framework for Cyclic Extension}

To bridge this gap, we draw on the theoretical foundations of cyclic
deformation developed for strike-slip systems (Nur and Mavko, 1974; Savage
and Prescott, 1978; Spence and Turcotte, 1979; Cohen and Kramer, 1984). These
works show that what appears as ``steady'' inter-rifting deformation is the
emergent superposition of visco-elastic responses from an infinite history of
discrete events. We model the lithosphere as an elastic crust overlying a
viscous substratum (e.g., Hetland, 2006), and derive an analytical solution
for the cyclic accumulation of displacement. The approach leverages the
mathematics of fractional calculus (Samko and Ross, 1993; Metzler and Klafter,
2000) to naturally reconcile post-diking transients with long-term plate
divergence, without requiring unphysical boundary conditions or an ad hoc
magmatic fraction of extension.

The remainder of this document presents the full derivations in two parts.
Appendix~\ref{sec:viscous} develops the purely viscous (Newtonian) model,
progressing from the governing equations through the step response of a single
dike to the Cyclo-Steady State (CSS) solution for an infinite sequence of
events. Appendix~\ref{sec:maxwell} extends the model to a Maxwell visco-elastic
asthenosphere and analyzes the resulting two-timescale behavior.

% ====================================================================
\appendix

\section{Elastic Crust over a Newtonian Viscous Half-Space}
\label{sec:viscous}

\subsection{Physical Setup and Governing Equations}

We consider the lithosphere as a coupled mechanical system consisting of an
elastic crust of thickness $H$ overlying a Newtonian viscous half-space
(asthenosphere) of viscosity $\eta$. The crust is modeled as a thin elastic
plate under plane strain conditions. The horizontal coordinate $x$ is
perpendicular to the rift axis; the vertical coordinate $z$ is oriented
upward, with the free surface at $z=H$ and the crust--mantle interface at
$z=0$.

Under the assumption of negligible vertical stress ($\sigma_{zz}\approx 0$),
the depth-averaged horizontal stress $\bar{\sigma}_{xx}$ is related to the
horizontal displacement $u(x,t)$ through Hooke's law:
\begin{equation}
\bar{\sigma}_{xx} = \frac{E}{1-\nu^2}\frac{\partial u}{\partial x},
\label{A1}
\end{equation}
where $E$ is Young's modulus and $\nu$ is Poisson's ratio. The factor
$(1-\nu^2)^{-1}$ accounts for the lateral confinement inherent in plane
strain. Mechanical equilibrium of the plate requires that the horizontal
gradient of this depth-averaged stress be balanced by the shear traction
exerted at its base by the underlying viscous mantle:
\begin{equation}
H\frac{\partial\bar{\sigma}_{xx}}{\partial x} + \sigma_{xz}^{\text{BOT}} = 0.
\label{A2}
\end{equation}

For a Newtonian viscous half-space, the basal shear stress is related to the
surface velocity field through a non-local convolution integral, the Stokes
kernel (Savage and Prescott, 1978):
\begin{equation}
\sigma_{xz}^{\text{BOT}}(x,t) =
  \int_{-\infty}^{+\infty}\frac{2\eta}{\pi(x-x')^2}\,
  \frac{\partial u}{\partial t}(x',t)\,dx',
\label{A3}
\end{equation}
interpreted as a Cauchy principal value. This non-local coupling is
particularly relevant in magmatically active rifts where the lower crust is
weakened by high heat flow and melt fractions (Holtzman, 2016; Plasman et al.,
2017). Substituting (\ref{A1}) into (\ref{A2}) and using (\ref{A3}) yields the
integro-differential equation governing the displacement field:
\begin{equation}
\frac{EH}{1-\nu^2}\frac{\partial^2 u}{\partial x^2}
+ \int_{-\infty}^{+\infty}\frac{2\eta}{\pi(x-x')^2}
  \frac{\partial u}{\partial t}(x',t)\,dx' = 0.
\label{A4}
\end{equation}

\subsection{Fourier Reduction to Fractional Diffusion}

We define the Fourier transform pair
\begin{equation*}
\hat{u}(k,t) = \int_{-\infty}^{+\infty}u(x,t)e^{-ikx}\,dx,\qquad
u(x,t) = \frac{1}{2\pi}\int_{-\infty}^{+\infty}\hat{u}(k,t)e^{ikx}\,dk.
\end{equation*}
The Fourier transform of $\partial^2 u/\partial x^2$ is $-k^2\hat{u}(k,t)$.
The transform of the Stokes kernel convolution is the classical result
(Savage and Prescott, 1978):
\begin{equation}
\mathcal{F}\!\left\{
  \int_{-\infty}^{+\infty}\frac{2\eta}{\pi(x-x')^2}
  \frac{\partial u}{\partial t}(x',t)\,dx'
\right\} = -2\eta|k|\frac{\partial\hat{u}}{\partial t}(k,t).
\label{A5}
\end{equation}
Applying the Fourier transform to (\ref{A4}) and using (\ref{A5}):
\begin{equation}
-\frac{EH}{1-\nu^2}k^2\hat{u} - 2\eta|k|\frac{\partial\hat{u}}{\partial t}=0.
\label{A6}
\end{equation}
Rearranging and introducing the \textbf{characteristic velocity}
\begin{equation}
\boxed{a := \frac{EH}{2\eta(1-\nu^2)}},
\label{A9}
\end{equation}
the equation simplifies to
\begin{equation}
\frac{\partial\hat{u}}{\partial t} = -a|k|\hat{u}.
\label{A10}
\end{equation}
In physical space the operator $|k|$ corresponds to the half-Laplacian
$(-\Delta)^{1/2}$, so (\ref{A10}) is equivalent to the fractional diffusion
equation
\begin{equation}
\frac{\partial u}{\partial t} = -a\,(-\Delta)^{1/2}u.
\label{A11}
\end{equation}
The parameter $a$ encodes the elastic--viscous coupling of the system: large
$a$ (stiff crust, low viscosity) implies rapid redistribution of post-diking
strain; small $a$ implies a broad, slowly evolving deformation pattern.
Equation (\ref{A10}) is a first-order ODE for each wavenumber $k$, with
solution
\begin{equation}
\hat{u}(k,t) = \hat{u}(k,0)\,e^{-a|k|t}.
\label{A12}
\end{equation}

\subsection{Step Response to a Single Dike Intrusion}
\label{sec:step}

We now solve (\ref{A12}) for the displacement field produced by a single
instantaneous dike intrusion at $t=0$. A symmetric dike opening by a total
amount $2\omega$ displaces the right side of the rift by $+\omega$ and the
left side by $-\omega$, with zero displacement elsewhere. Exploiting the
antisymmetry of the problem (working for $x>0$), the initial condition is
represented by a Heaviside function:
\begin{equation}
u(x,t=0^+) = \omega\,\mathcal{H}(-x).
\label{A13}
\end{equation}

\paragraph{Fourier transform of the initial condition.}
Applying the Fourier transform to (\ref{A13}):
\begin{equation}
\hat{u}(k,0) = \omega\infint \mathcal{H}(-x)\,e^{-ikx}\,dx
             = \omega\int_{-\infty}^{0}e^{-ikx}\,dx.
\label{A14}
\end{equation}
For real $k$ this integral is treated in the distributional sense by adding a
convergence regulator $\varepsilon>0$:
\begin{equation*}
\int_{-\infty}^{0}e^{-ikx}\,dx
= \lim_{\varepsilon\to 0}\int_{-\infty}^{0}e^{(-ik+\varepsilon)x}\,dx
= \lim_{\varepsilon\to 0}\frac{1}{\varepsilon - ik}.
\end{equation*}
By the Sokhotski--Plemelj theorem, the limit gives
\begin{equation*}
\int_{-\infty}^{0}e^{-ikx}\,dx = \pi\delta(k) + \frac{i}{k},
\end{equation*}
so the initial spectrum is
\begin{equation}
\hat{u}(k,0)= \omega\pi\delta(k) + \frac{i\omega}{k}.
\label{A15}
\end{equation}

\paragraph{Evolved solution and inverse transform.}
Substituting (\ref{A15}) into (\ref{A12}), and noting that the $\delta(k)$
term is only active at $k=0$ (where the exponential equals 1):
\begin{equation}
\hat{u}(k,t) = \omega\pi\delta(k) + \frac{i\omega}{k}\,e^{-a|k|t}.
\label{A16}
\end{equation}
Taking the inverse Fourier transform:
\begin{equation}
u(x,t) = \frac{1}{2\pi}\infint
  \left(\omega\pi\delta(k) + \frac{i\omega}{k}\,e^{-a|k|t}\right)e^{ikx}\,dk.
\label{A17}
\end{equation}
We split this into $I_1$ and $I_2$. The delta term gives a constant background:
\begin{equation}
I_1 = \frac{\omega}{2}.
\label{A18}
\end{equation}
For $I_2$, using Euler's identity and noting that $1/k\cdot\cos(kx)$ is odd
while $1/k\cdot\sin(kx)$ is even:
\begin{equation}
I_2 = -\frac{\omega}{\pi}\int_0^{\infty}\frac{\sin(kx)}{k}\,e^{-akt}\,dk.
\label{A19}
\end{equation}
This is a standard Dirichlet-type integral with value
$-\tfrac{\omega}{\pi}\arctan\!\left(\tfrac{x}{at}\right)$.
Combining $I_1$ and $I_2$ and applying the complementary arctan identity
$\arctan(z)+\arctan(1/z)=\pi/2$, the \textbf{step displacement} is:
\begin{equation}
\boxed{u^{\text{STEP}}(x,t) = \frac{\omega}{\pi}\arctan\!\left(\frac{at}{x}\right).}
\label{A22}
\end{equation}
Differentiating with respect to $t$ yields the \textbf{step velocity},
a Lorentzian profile that broadens in time:
\begin{equation}
v^{\text{STEP}}(x,t) = \frac{\omega a}{\pi}\,\frac{x}{x^2 + a^2t^2}.
\label{A23}
\end{equation}
Both profiles satisfy the correct limits: $u^{\text{STEP}}\to\omega/2$ and
$v^{\text{STEP}}\to 0$ as $t\to\infty$, consistent with the far-field plate
divergence.

\subsection{Cyclic Accumulation and the Cyclo-Steady State}

\subsubsection{Finite sum over past intrusions}

For a rift that has experienced $n = \lfloor t/T\rfloor$ dike intrusions at
intervals of $T$ years, the total displacement is the superposition of
individual step responses:
\begin{equation}
u(x,t) = \frac{\omega}{\pi}\sum_{j=1}^{n}
  \arctan\!\left(\frac{a(t-jT)}{x}\right).
\label{A24}
\end{equation}
Differentiating with respect to $t$, the velocity profile is:
\begin{equation}
v(x,t) = \frac{\omega a}{\pi}\sum_{j=1}^{\lfloor t/T\rfloor}
  \frac{x}{x^2 + a^2(t-jT)^2}.
\label{A25v}
\end{equation}

\subsubsection{The CSS limit and Digamma representation}

After many cycles the system enters a periodic state. Introducing
$\tau\in[0,T)$ as the time elapsed since the most recent intrusion,
the \textbf{Cyclo-Steady State} (CSS) velocity is obtained by extending the
sum in (\ref{A25v}) to an infinite history:
\begin{equation}
v_{\text{CSS}}(x,\tau) = \frac{\omega a}{\pi}
  \sum_{j=1}^{\infty}\frac{x}{x^2 + a^2(\tau-jT)^2}.
\label{A25}
\end{equation}
Because this series converges, it can be expressed in closed form using the
Digamma function $\psi(z)$. Using the partial-fraction identity
\begin{equation*}
\frac{x}{x^2+a^2(\tau-jT)^2}
= \frac{1}{2i}\left[
    \frac{1}{a(\tau-jT)-ix} - \frac{1}{a(\tau-jT)+ix}
  \right],
\end{equation*}
factoring $-aT$ from each term, and substituting into (\ref{A25}), the
double sum over $j$ takes the form of differences of Digamma functions.
Recalling the series representation $\psi(z+1) = -\gamma + \sum_{n=1}^\infty
\!\left(1/n - 1/(n+z)\right)$, the Euler--Mascheroni constant $\gamma$ and
the $1/n$ terms cancel when the difference of two conjugate Digamma functions
is taken:
\begin{equation}
\psi(Z)-\psi(\bar{Z}) = 2i\,\IM\{\psi(Z)\}.
\end{equation}
With $Z = \tau/T + ix/(aT)$, the CSS velocity collapses to:
\begin{equation}
\boxed{v_{\text{CSS}}(x,\tau) =
  \frac{2\omega}{\pi T}\,\IM\!\left\{\psi\!\left(
    \frac{\tau}{T}+i\frac{x}{aT}\right)\right\}.}
\label{A32}
\end{equation}

\subsubsection{Displacement and mass conservation}

Integrating (\ref{A32}) over $\tau$ using the antiderivative
$d[\ln\Gamma(z)]/dz = \psi(z)$, the \textbf{CSS displacement} is:
\begin{equation}
\boxed{u_{\text{CSS}}(x,\tau) =
  \frac{2\omega}{\pi}\,\IM\!\left\{
    \ln\Gamma\!\left(\frac{\tau}{T}+i\frac{x}{aT}\right)
    - \ln\Gamma\!\left(i\frac{x}{aT}\right)
  \right\}.}
\label{A33}
\end{equation}
A key property of the CSS solution is that, over one complete cycle, the
cumulative displacement at any fixed $x>0$ equals $\omega$:
\begin{equation}
\int_0^T v_{\text{CSS}}(x,\tau)\,d\tau = \omega, \qquad \forall\,x>0.
\label{A34}
\end{equation}
This means the entire right half of the plate advances by $\omega$ per cycle
regardless of distance from the rift, exactly consistent with far-field
kinematics prescribed by global plate motion models (Saria et al., 2014;
Stamps et al., 2021). The CSS thus simultaneously reconciles high-frequency
magmatic pulses with long-term secular divergence, without any unphysical
boundary conditions.

% ====================================================================
\section{Maxwell Visco-Elastic Asthenosphere}
\label{sec:maxwell}

\subsection{Constitutive Model and the Correspondence Principle}

In Appendix~\ref{sec:viscous} the asthenosphere was modeled as a purely
Newtonian fluid. While adequate for long-term tectonic processes, this
assumption fails to capture the immediate elastic response of the mantle to a
sudden dike intrusion. We therefore upgrade the asthenosphere to a
\textbf{Maxwell body}: a Hookean spring and a Newtonian dashpot connected in
series, characterized by shear modulus $\mu_m$ and viscosity $\eta$. The 1D
constitutive relation for the deviatoric stress $\tau$ and strain $\varepsilon$
is:
\begin{equation}
\dot{\varepsilon} = \frac{\dot{\tau}}{2\mu_m} + \frac{\tau}{2\eta}.
\label{B1}
\end{equation}
The \textbf{Maxwell relaxation time} $\tau_R := \eta/\mu_m$ characterizes the
timescale over which the material transitions from elastic to viscous behavior.

To adapt the Stokes kernel (\ref{A3}) for a visco-elastic half-space, we
employ the \textbf{Correspondence Principle} of linear visco-elasticity: the
solution to a visco-elastic problem is obtained from the corresponding elastic
solution by replacing static moduli with their complex, frequency-dependent
counterparts in the Laplace domain. Denoting the Laplace transform by
$\tilde{f}(s) = \mathcal{L}\{f(t)\}$ and applying it to (\ref{B1}) with zero
initial conditions:
\begin{equation}
s\tilde{\varepsilon} = \tilde{\tau}\!\left(\frac{s}{2\mu_m}+\frac{1}{2\eta}\right).
\label{B2}
\end{equation}
Solving for stress gives $\tilde{\tau} = 2\tilde{\mu}(s)\tilde{\varepsilon}$,
where the effective, frequency-dependent shear modulus of the Maxwell body is:
\begin{equation}
\tilde{\mu}(s) = \frac{\mu_m\eta\,s}{\mu_m + \eta s}
               = \mu_m\frac{\tau_R s}{1+\tau_R s}.
\label{B4}
\end{equation}

\subsection{The Visco-Elastic Stokes Kernel}

For a purely elastic half-space the basal traction relates to surface
displacement via a kernel proportional to $\mu_m$:
\begin{equation}
\sigma_{xz}^{\text{BOT, el}}(x) =
  \int_{-\infty}^{+\infty}\frac{2\mu_m}{\pi(x-x')^2}\,u(x')\,dx'.
\label{B5}
\end{equation}
Applying the Correspondence Principle---replacing $\mu_m$ with $\tilde{\mu}(s)$
in the Laplace domain---and factoring out $\eta = \mu_m\tau_R$:
\begin{equation}
(1+\tau_R s)\,\tilde{\sigma}_{xz}^{\text{BOT}}(x,s) =
  \int_{-\infty}^{+\infty}\frac{2\eta}{\pi(x-x')^2}
  \left[s\tilde{u}(x',s)\right]dx'.
\label{B8}
\end{equation}
Inverting the Laplace transform (multiplication by $s$ corresponds to
$\partial/\partial t$) yields the time-domain visco-elastic Stokes kernel as a
differential relation governing the basal traction:
\begin{equation}
\sigma_{xz}^{\text{BOT}} + \tau_R\frac{\partial\sigma_{xz}^{\text{BOT}}}{\partial t}
= \int_{-\infty}^{+\infty}\frac{2\eta}{\pi(x-x')^2}
  \frac{\partial u(x',t)}{\partial t}\,dx'.
\label{B9}
\end{equation}
In the limit $\tau_R\to 0$ the left-hand side reduces to
$\sigma_{xz}^{\text{BOT}}$ alone, recovering the purely viscous Stokes kernel
(\ref{A3}).

\subsection{Governing Equation in Fourier Space}

The equilibrium equation for the elastic crust (\ref{A2}) gives:
\begin{equation}
\sigma_{xz}^{\text{BOT}} = -\frac{EH}{1-\nu^2}\frac{\partial^2 u}{\partial x^2}.
\label{B10}
\end{equation}
Substituting (\ref{B10}) and its time derivative into (\ref{B9}), and
then applying the spatial Fourier transform using (\ref{A5}):
\begin{equation}
\frac{EH}{1-\nu^2}k^2\hat{u}
+ \tau_R\frac{EH}{1-\nu^2}k^2\frac{\partial\hat{u}}{\partial t}
= -2\eta|k|\frac{\partial\hat{u}}{\partial t}.
\label{B13}
\end{equation}
Dividing by $2\eta$, using the definition of $a$ from (\ref{A9}), and
grouping $\partial\hat{u}/\partial t$ terms:
\begin{equation}
\frac{\partial\hat{u}}{\partial t}
= -\left(\frac{a|k|}{1+\tau_R a|k|}\right)\hat{u}.
\label{B16}
\end{equation}
This is the visco-elastic analogue of (\ref{A10}). The denominator
$1+\tau_R a|k|$ acts as a low-pass filter on the relaxation rate, slowing the
decay of high-wavenumber (short-wavelength, near-field) features relative to
the purely viscous case.

\subsection{Step Response and Asymptotic Behavior}

\subsubsection{Solution for an instantaneous dike intrusion}

Equation (\ref{B16}) is separable in $k$, giving a decaying exponential for
each wavenumber:
\begin{equation}
\hat{u}(k,t) = \hat{u}(k,0)\exp\!\left(-\frac{a|k|}{1+\tau_R a|k|}\,t\right).
\label{B17}
\end{equation}
Applying the same Heaviside initial condition (\ref{A15}) and noting that the
$\delta(k)$ term evaluates trivially:
\begin{equation}
\hat{u}(k,t) = \omega\pi\delta(k)
  + \frac{i\omega}{k}\exp\!\left(-\frac{a|k|}{1+\tau_R a|k|}\,t\right).
\label{B20}
\end{equation}
Taking the inverse Fourier transform and following the same parity argument
as in Section~\ref{sec:step}:
\begin{equation}
u(x,t) = \frac{\omega}{2}
  - \frac{\omega}{\pi}\int_0^{+\infty}
    \frac{\sin(kx)}{k}\exp\!\left(-\frac{akt}{1+\tau_R ak}\right)dk.
\label{B22}
\end{equation}
Unlike the purely viscous case, the rational exponent prevents a closed-form
evaluation of this integral, which must be computed numerically. The surface
velocity follows by differentiating (\ref{B22}) with respect to $t$:
\begin{equation}
v(x,t) = \frac{\omega a}{\pi}\int_0^{+\infty}
  \frac{\sin(kx)}{1+\tau_R ak}
  \exp\!\left(-\frac{akt}{1+\tau_R ak}\right)dk.
\label{B23v}
\end{equation}

\subsubsection{Asymptotic limits}

The structure of the exponent in (\ref{B22}) reveals two physically distinct
regimes that illuminate the two-timescale behavior of the Maxwell model.

\paragraph{Long-wavelength / long-time limit ($k\to 0$ or $t\to\infty$).}
For broad spatial features, $1+\tau_R ak\approx 1$ and the exponential
reduces to $\exp(-akt)$, recovering the purely viscous Dirichlet integral of
Section~\ref{sec:step}. Over long timescales or large distances, the elastic
memory of the asthenosphere fades and the model converges to the Newtonian
viscous limit.

\paragraph{Short-wavelength / short-time limit ($k\to\infty$ or $t\to 0$).}
For features localized near the rift axis, $1+\tau_R ak\approx\tau_R ak$,
and the exponent saturates:
\begin{equation}
\lim_{k\to\infty}\left(-\frac{akt}{1+\tau_R ak}\right) = -\frac{t}{\tau_R}.
\label{B24}
\end{equation}
At early times $t\ll\tau_R$ the exponential approaches unity, so the integral
evaluates to $\pi/2$ and no near-field deformation occurs outside the rift
interface. For $t\sim\tau_R$, near-field deformation decays exponentially at
the fixed rate $1/\tau_R$, representing pure Maxwell relaxation of the mantle
independent of the elastic crust. Together, (\ref{B23v}) and (\ref{B24})
show that the Maxwell model introduces a short-time, near-field stiffening
relative to the viscous model, before both converge to the same CSS at
$t\gg\tau_R$.

\bibliographystyle{alpha}
\end{document}
